\chapter{Introduction}
\label{ch:introduction}

Causal discovery studies how causal structure can be inferred, tested or refined from data and background knowledge. This is important because causal structure supports explanation, prediction under intervention and decision-making in settings where direct experimentation may be infeasible, unethical or prohibitively expensive \cite{glymour2019review,vowels2022d}. A common approach is to represent causal hypotheses as directed acyclic graphs, where directed edges encode possible direct causal relations. Under assumptions such as the causal Markov condition and faithfulness, graphical separation relations can be connected to statistical conditional independences. This makes it possible to use observational data to constrain the class of causal graphs compatible with the evidence.

However, observational causal discovery is difficult. Conditional-independence tests are finite-sample statistical procedures, and their outputs may vary with sample size, test choice, significance threshold and modelling assumptions. More generally, causal-discovery algorithms differ in how they convert uncertain statistical evidence into graph structure, whether through constraint-based testing, score-based search, functional causal assumptions or continuous optimisation \cite{glymour2019review}. In small or noisy settings, the resulting evidence may be incomplete, unstable or mutually inconsistent. This motivates a central problem for this project: how can causal-discovery evidence be represented and reasoned with when it is defeasible rather than certain?

Formal argumentation provides one possible answer. In argumentation, conclusions can be supported by arguments, challenged by counterarguments and accepted only under a specified semantics. Assumption-Based Argumentation (ABA) is particularly relevant because it represents reasoning in terms of rules, defeasible assumptions and contraries \cite{bondarenko1997abstract,toni2014tutorial}. A causal claim can therefore be treated not as an irrevocable fact, but as an assumption that may be supported or attacked by other information. This is the motivation behind Causal ABA, which represents causal-discovery evidence inside an ABA framework. In Causal ABA, assumptions can encode candidate arrows, no-edge claims and conditional-independence information, while rules encode graph-theoretic constraints such as acyclicity and d-separation. Stable extensions can then be interpreted as compatible causal graph hypotheses \cite{russo2024argumentative}.

A separate but closely related line of work is ABA Learning. Whereas Causal ABA represents and evaluates causal-discovery evidence in an ABA framework, ABA Learning studies how an ABA framework can itself be transformed or extended from background knowledge and positive and negative examples. The learning process uses transformation rules such as Rote Learning, Folding, Assumption Introduction and Subsumption to construct a framework in which positive examples are accepted and negative examples are not accepted under a chosen semantics \cite{proietti2022learning,de2023aba,de2024learning}. Recent work on Greedy ABA Learning further shows that the space of possible transformations can be large and that controlling the learning process is an important practical issue \cite{de2025greedy}.

The gap addressed by this project lies between these two strands. Causal ABA gives an argumentative representation of causal-discovery evidence, but it does not by itself study how the relevant ABA structures should be learnt from data. ABA Learning gives machinery for inducing ABA frameworks from examples, but existing ABA Learning methods are not specialised for causal graphs, conditional-independence evidence, d-separation or acyclicity constraints. The project therefore investigates the interaction between causal discovery and ABA Learning. This project pursues the integration direction in which Causal ABA guides ABA Learning. It investigates whether causal information such as candidate arrows, no-edge claims, conditional-independence evidence, acyclicity and d-separation can be used as argumentative background knowledge to constrain, prioritise or interpret ABA Learning transformations. The ultimate aim is to learn causal relationships from data in an argumentative fashion.

The project is deliberately scoped. The current implementation used in the interim experiments is not a full implementation of Russo et al.'s Causal ABA framework. It does not yet implement the full \(\mathsf{arr}\), \(\mathsf{noe}\) and independence-assumption machinery, nor the d-separation encoding and stable-extension-as-DAG construction used in Causal ABA. Instead, the current empirical work studies a narrower target-wise ABA Learning pipeline: a tabular dataset is converted into ABA background predicates, ABA Learning is run for a selected target variable, and the variables appearing in learnt target-rule bodies are compared with the true direct parents in a known synthetic data-generating graph. This parent-set comparison is a diagnostic proxy, not a proof of causal discovery. The initial experiments in this report therefore serve as groundwork: they characterise what unguided ABA Learning does on small causal motifs before stronger Causal ABA-inspired guidance is introduced.

The research questions for the project are:

\begin{enumerate}
    \item \textbf{RQ1:} When does unguided ABA Learning recover parent variables in learnt target rules, rather than ancestors, descendants, siblings, correlated non-parents or representation artefacts?
    \item \textbf{RQ2:} Can Causal ABA-style evidence and graph constraints be used to guide ABA Learning towards rules or frameworks that are more causally meaningful, interpretable, robust or efficient?
    \item \textbf{RQ3:} How does the resulting causally guided ABA Learning bridge compare with representative causal-discovery methods on controlled canonical examples?
\end{enumerate}

At this interim stage, the main empirical contribution concerns RQ1. The experiments investigate chain, fork and collider motifs under binary, categorical and discretised continuous representations. They show that the current target-wise ABA Learning pipeline can run end-to-end and can recover true parent sets in selected idealised cases. They also show important limitations: recovery is sensitive to motif, encoding and noise; collider-like two-parent mechanisms are difficult for the current configuration; and noisy scaled settings can produce no-solution or error outcomes even when the computation terminates. These findings do not invalidate the project aim. After conducting a more principled and systematic investigation to establish how and when unguided ABA Learning can recover causality or its proxies, the project will transition to its core exploratory work of using Causal ABA-style causal information to guide ABA Learning.

The remainder of this report develops this argument in four stages. The Literature Review positions the project within causal discovery, finite-sample uncertainty, argumentation-based causal discovery and ABA Learning, and identifies the integration gap validating a fusion of Causal ABA and ABA Learning. The Background chapter fixes the formal notions needed later, including DAGs and conditional independence, ABA frameworks and stable extensions, Causal ABA, and ABA Learning transformation rules. The Experimentation chapter reports the initial parent-set recovery investigation and leverages the its intepretation and limitations to motivate the Project Plan. The plan sets out the remaining work: closing the parent-set diagnostic, designing a Causal ABA-guided learning bridge, evaluating guided and unguided learning on controlled examples, and comparing the resulting approach with representative causal-discovery baselines.